% Math
\usepackage{amsmath}
\usepackage{upgreek}

% Miscelaneous
\usepackage{pgfpages}
\usepackage{csquotes}
\usepackage{placeins}
\usepackage[onehalfspacing]{setspace}

% Captions
\usepackage[labelformat=empty]{caption}

% Colors
\usepackage{xcolor}

% Requirements
\RequirePackage{tikz}
\usetikzlibrary{calc}
\usetikzlibrary{arrows,shapes,positioning}

% For \smiley{} and \frowney{}
\usepackage{wasysym}

% Pseudocode
\usepackage[vlined, linesnumbered, ruled]{algorithm2e}
\SetKwComment{Comment}{/* }{*/}
\SetKwInput{KwData}{Input}
\SetKwInput{KwResult}{Output}
\IncMargin{1.5em}

% nice icons
%\usepackage{fontawesome}

% IMPORTANT NOTE: outputdir needs to be equal to output/aux dir of xelatex
%\usepackage[outputdir=output]{minted}
\usepackage{appendixnumberbeamer}

% Tables
\usepackage{booktabs}
\usepackage{multicol}
\usepackage{multirow}
\usepackage{colortbl}

% Date formatting
\usepackage[short,nodayofweek]{datetime}

% Für Vektoren
\newcommand{\myvec}[1]{\ensuremath{\begin{pmatrix}#1\end{pmatrix}}}
\def\spvec#1{\left(\vcenter{\halign{\hfil$##$\hfil\cr \spvecA#1;;}}\right)}
\def\spvecA#1;{\if;#1;\else #1\cr \expandafter \spvecA \fi}

% Hyperlinks
\definecolor{citeblue}{rgb}{0,0.3,0.55}
\RequirePackage{hyperref}
\hypersetup{colorlinks,linkcolor={lightgray},citecolor={citeblue},urlcolor={citeblue}}
\let\oldciteauthor=\citeauthor
\def\citeauthor#1{\hypersetup{citecolor=white}\oldciteauthor{#1}}

% Added
\usepackage{circuitikz}
\usepackage{grfext}


% Custom Makros
\newcommand{\bangle}[1]{\langle #1 \rangle}
\newcommand{\semsof}[1]{\ensuremath{\llbracket #1 \rrbracket}}
\newcommand{\impl}{\Rightarrow}
\newcommand{\reason}[1]{\qquad \mid #1}
\newcommand{\edge}{\ensuremath{(e,e')}}
\newcommand{\edgein}{\ensuremath{(e,e') \in }}
\newcommand{\egdein}{\ensuremath{(e',e) \in }}
\newcommand{\defof}{\text{Def. of }~}
\newcommand{\inv}{\ensuremath{{-1}}}
\newcommand{\prf}{\mathcal{P}}
\newcommand{\catprf}{\ensuremath{\mathcal{P}^{CAT}}}
\newcommand{\I}{\mathcal{I}}

% Reelle, Natürliche, Ganze, Rationale Zahlen
\newcommand{\R}{\ensuremath{\mathbb{R}}}
\newcommand{\N}{\ensuremath{\mathbb{N}}}
\newcommand{\Z}{\ensuremath{\mathbb{Z}}}
\newcommand{\Q}{\ensuremath{\mathbb{Q}}}
\newcommand{\E}{\ensuremath{\mathbb{E}}}
\newcommand{\B}{\ensuremath{\mathfrak{B}}}
\newcommand{\C}{\ensuremath{\mathfrak{C}}}


% Kürzel für Basisrelationen
\definecolor{rfgreen}{rgb}{0.2,0.55,0.2}
\newcommand{\0}{\textsf{\upshape 0}} % \upshape ensures that the relations remain their look in lemmata and inferrules
\newcommand{\po}{\textsf{\upshape po}}
\newcommand{\co}{\textsf{\textcolor{orange}{\upshape co}}}
\newcommand{\rf}{\textsf{\textcolor{rfgreen}{\upshape rf}}}
\newcommand{\ext}{\textsf{\upshape ext}}
\newcommand{\rmw}{\textsf{\upshape rmw}}
\newcommand{\hb}{\textsf{\textcolor{purple}{\upshape hb}}}
\newcommand{\loc}{\textsf{\upshape loc}}
\newcommand{\fr}{\textsf{\textcolor{violet}{\upshape fr}}}
\newcommand{\id}{\textsf{\upshape id}}
\newcommand{\true}{\textit{true}}
\newcommand{\false}{\textit{false}}
\newcommand{\rel}{\bangle{rel}}

%\newcommand{\rel}[1]{\textsf{#1}}
\newcommand{\relI}{\textsf{\upshape I}}
\newcommand{\relW}{\textsf{\upshape [W]}}
\newcommand{\relR}{\textsf{\upshape [R]}}
\newcommand{\relr}{\textsf{\upshape r}}
\newcommand{\Relr}{\relr}
\newcommand{\rels}{\textsf{\upshape s}}
\newcommand{\Rels}{\rels}
\newcommand{\relt}{\textsf{\upshape t}}


% Definitionsboxen
\usepackage[most]{tcolorbox}
%\definecolor{darkgray}{RGB}{70,90,70}
\definecolor{examblue}{rgb}{0,0.2,0.55}
\newenvironment{defi}[1]{%
\begin{tcolorbox}[blanker, before skip=5pt, after skip = 5pt, colback=darkgray!5!white, colframe=darkgray!75!black, title={{\sffamily} \textbf{#1}}]
}{
\end{tcolorbox}}


\usetikzlibrary{decorations.pathreplacing,calc}
\newcommand{\tikzmark}[1]{\tikz[overlay,remember picture] \node (#1) {};}

\newcommand*{\AddNote}[5]{%
	\begin{tikzpicture}[overlay, remember picture]
		\draw [decoration={brace,amplitude=0.5em},decorate,ultra thick,color=#4]
		($(#3)!(#2.south)!($(#3)-(0,1)$)$) --
		($(#3)!(#1.north)!($(#3)-(0,1)$)$)
		node [align=center, text width=0.5cm, pos=0.5, anchor=east] {#5};
	\end{tikzpicture}
}%

\usepackage{algpseudocode}% http://ctan.org/pkg/algorithmicx
%\usepackage{algorithm}% http://ctan.org/pkg/algorithm

% captions for subfigures.
\usepackage{subcaption}


% Clever references.
\usepackage{hyperref}
\usepackage{cleveref}
\usepackage{amssymb}